\item Dos núcleos con números atómicos $Z_1$ y $Z_2$ se aproximan mutuamente con una energía total $E$.
\begin{enumerate}
    \item Si interactúan culombianamente hasta una distancia crítica de $1.00 \times 10^{-14}\text{ m}$ donde la fuerza nuclear atractiva toma el control, encuentre la energía mínima $E$ en función de $Z_1$ y $Z_2$ requerida para fusionarse.
    \item Indique cómo depende cuantitativamente esta energía de los números atómicos.
    \item Si $Z_1 + Z_2 = 60$, demuestre si es energéticamente más favorable para la fusión termonuclear elegir la combinación $Z_1 = 1$, $Z_2 = 59$ o $Z_1 = Z_2 = 30$.
    \item Evalúe numéricamente la energía mínima clásica de fusión para las reacciones termonucleares $\text{D}-\text{D}$ y $\text{D}-\text{T}$.
\end{enumerate}